\chapter{Analyse et conception}
\label{chap:02-analyse-conception}
\markboth{Analyse et conception}{ }
\clearpage

\section{De la problématique aux exigences}
\label{sec:02-exigences}

Le besoin initial est exprimé en termes de suivi de demandes. Pour devenir développable, il est transformé en exigences observables : qui réalise l'action, sur quelle donnée, à quel moment et sous quelles conditions ? Les exigences fonctionnelles couvrent l'authentification, la gestion des demandes et clients, le workflow, l'administration et l'assistance IA. Les exigences non fonctionnelles ajoutent des contraintes de sécurité, cohérence transactionnelle, maintenabilité, ergonomie, accessibilité et traçabilité.

\begin{table}[htbp]
\centering
\caption{Synthèse des exigences principales}
\label{tab:02-exigences}
\begin{tabular}{p{0.27\textwidth}p{0.62\textwidth}}
\textbf{Domaine} & \textbf{Attente principale}\\
Authentification & Identifier un utilisateur actif et transmettre un JWT aux requêtes protégées.\\
Demandes & Créer, consulter, filtrer, affecter, traiter, résoudre, clôturer ou annuler une demande.\\
Administration & Gérer les comptes, rôles métier et états actifs.\\
Sécurité & Appliquer RBAC et contrôles contextuels côté backend.\\
IA & Proposer une catégorie et une priorité sans application automatique.\\
\end{tabular}
\end{table}

\section{Règles métier et cycle de vie}
\label{sec:02-regles}

Une demande ne se résume pas à un statut modifiable librement. Chaque transition correspond à une action métier explicite. L'affectation associe une demande à un agent ; le traitement permet de renseigner les éléments nécessaires ; la résolution est déclarée par l'Agent puis examinée par le Responsable technique. Le refus renvoie la demande vers une reprise. L'annulation exige un motif et reste distincte de la clôture. Les changements métier et leur historisation sont atomiques.

Le diagramme d'activité répond à la question des transitions possibles depuis chaque état. Le diagramme de cas d'utilisation montre quel acteur demande quelle capacité. Les séquences validées détaillent les scénarios majeurs sans exposer le détail des contrôleurs et repositories ; le lecteur comprend d'abord le comportement, avant l'implémentation.

\section{Architecture retenue}
\label{sec:02-architecture}

La solution combine Angular, une API REST Spring Boot et PostgreSQL. Le backend est un monolithe modulaire : les fonctionnalités partagent un déploiement, mais sont séparées autour de `request`, `identity / administration`, `security` et `ai`. Cette organisation répond au besoin de simplicité opérationnelle tout en préparant une évolution maîtrisée.

La séparation interne distingue présentation, application, domaine et persistance. Le Controller traduit HTTP vers un cas d'utilisation ; le Service orchestre ce cas et sa transaction ; le domaine porte les invariants ; le Repository encapsule la persistance ; les DTO stabilisent les échanges et séparent contrat API et entité JPA.

\section{SOLID et patterns : analyse limitée aux preuves}
\label{sec:02-solid-patterns}

Le principe de responsabilité unique est observable dans la séparation entre contrôleurs, services, objets métier et repositories. Cette organisation améliore la lisibilité, mais ne prouve pas une conformité absolue à SOLID. Le principe d'inversion des dépendances est visible dans `AnalyseIaPort` : le cœur applicatif dépend d'une abstraction et l'adaptateur Groq fournit l'intégration concrète.

Le projet n'utilise pas le State Pattern pour le cycle de vie. Les états sont portés par un enum et des opérations métier explicites, solution adaptée au nombre limité d'états de la V1. Cette décision pourra être réévaluée si le workflow devient beaucoup plus complexe.

\section{Conception des données}
\label{sec:02-donnees}

Les entités persistées sont `Utilisateur`, `Client`, `DemandeTechnique` et `HistoriqueDemande`. La référence métier est générée côté serveur, unique et immuable ; l'email utilisateur est unique et normalisé. JPA et Hibernate assurent le mapping objet-relationnel, Spring Data JPA fournit les repositories et Flyway applique les migrations SQL versionnées. Les instants utilisent `Instant` côté Java et `timestamptz` côté PostgreSQL \parencite{postgresql-datetime-types}.

\section{API, sécurité et session Angular}
\label{sec:02-securite}

L'API REST expose des opérations orientées métier et retourne des DTO. Le backend est l'autorité pour l'identité, l'état actif, les rôles courants et les permissions. Le contrôle d'accès fondé sur les rôles (RBAC) est complété par des contrôles contextuels, notamment pour l'affectation et l'auto-désactivation.

Le JSON Web Token (JWT) est un format compact de claims standardisé par l'IETF \parencite{rfc7519-jwt}. Angular conserve la session en mémoire et dans `sessionStorage` sous `safyron.session.v1`. Il effectue des contrôles locaux de structure et d'expiration, sans remplacer la validation backend de la signature et des droits. Les mots de passe sont traités par `PasswordEncoder` et la configuration sensible est externalisée \parencite{spring-security-password-storage,spring-boot-externalized-configuration}.

\section{Conception de l'assistance IA}
\label{sec:02-conception-ia}

L'IA réduit l'effort de qualification, mais ne prend pas la décision. L'entrée est limitée au titre et à la description ; la sortie propose une catégorie et une priorité. `AnalyseIaPort` isole le cœur applicatif de l'adaptateur Groq. La proposition est temporaire, modifiable et appliquée après un clic explicite. En cas d'indisponibilité, la saisie manuelle reste possible.
